% Standalone source; no external bibliography or figures are required.
% Build without writing generated files into the repository:
% mkdir -p /tmp/sfcc-correctness-proofs
% latexmk -pdf -interaction=nonstopmode -halt-on-error \
%   -outdir=/tmp/sfcc-correctness-proofs \
%   docs/research/sfcc-correctness-proofs.tex
\documentclass[11pt]{article}
\usepackage[T1]{fontenc}
\usepackage[utf8]{inputenc}
\usepackage{lmodern}
\usepackage[margin=1in]{geometry}
\usepackage{amsmath,amssymb,amsthm,mathtools}
\usepackage{booktabs,array}
\usepackage{enumitem}
\usepackage{xurl}
\usepackage[colorlinks=true,linkcolor=blue,citecolor=blue,urlcolor=blue]{hyperref}
\hypersetup{pdftitle={SFCC correctness proofs: polyhedral completions and obstructions},
  pdfsubject={Exact rational polyhedral contouring; theory, not implementation certification}}
\setlist{itemsep=3pt,topsep=5pt}
\setlength{\emergencystretch}{2em}
\newtheorem{theorem}{Theorem}[section]
\newtheorem{lemma}[theorem]{Lemma}
\newtheorem{proposition}[theorem]{Proposition}
\newtheorem{corollary}[theorem]{Corollary}
\theoremstyle{definition}
\newtheorem{definition}[theorem]{Definition}
\newtheorem{algorithm}[theorem]{Algorithm}
\theoremstyle{remark}
\newtheorem{remark}[theorem]{Remark}
\newcommand{\R}{\mathbb R}
\newcommand{\Q}{\mathbb Q}
\newcommand{\Z}{\mathbb Z}
\newcommand{\Fq}{\mathbb F_q}
\newcommand{\cl}{\operatorname{cl}}
\newcommand{\intr}{\operatorname{int}}
\newcommand{\relint}{\operatorname{relint}}
\newcommand{\reg}{\operatorname{reg}}
\newcommand{\dist}{\operatorname{dist}}
\newcommand{\diam}{\operatorname{diam}}
\newcommand{\conv}{\operatorname{conv}}
\newcommand{\aff}{\operatorname{aff}}
\newcommand{\supp}{\operatorname{supp}}
\newcommand{\norm}[1]{\left\lVert #1\right\rVert}
\newcommand{\abs}[1]{\left\lvert #1\right\rvert}
\newcommand{\Binf}[2]{B_\infty(#1,#2)}
\newcommand{\Ninf}[2]{N_\infty(#1,#2)}
\newcommand{\symdiff}{\mathbin{\triangle}}

\title{SFCC Correctness Proofs:\\
  Polyhedral Completions, Refinement Obstructions,\\
  and the Remaining Curved-Surface Obligations}
\author{Galacticad research documentation}
\date{October 13, 2026}

\begin{document}
\maketitle

\begin{abstract}
Stratified Feature-Conforming Contouring (SFCC) proposes to compile surface
strata, feature curves, and corners; refine an octree; compute each shared face
trace once; and fill smooth, edge, and corner cells. Those design ingredients
alone do not prove feature completeness, termination, manifoldness, or an
intersection-free embedding.

This document gives four separate results for exact rational polyhedra.
First, a support-coordinate certificate makes the simple-cell construction
terminate for regularized Boolean combinations of axis-aligned boxes.
Second, an explicit rational tetrahedron prevents termination of the original
one-feature-or-one-corner rule on a fixed face-balanced dyadic octree.
Third, a finite arrangement completion gives an exact, unrestricted fallback.
Fourth, the main theorem proves termination and exactness for a
\emph{protected-neighborhood completion}: corner and edge volume cells are
frozen, arbitrary depth differences are allowed at their boundaries, complete
shared face graphs are constructed exactly, and the remaining cells are
certified planar. The resulting geometric simplicial complex equals the solid
boundary as a set and preserves every true corner and straight crease.

These are constructive mathematical specifications, not certificates for the
current exporter. Curved carriers, floating-point output, triangle quality, and
practical complexity are outside the proved theorem. Local residual and
conditioning estimates are supplied to state some of the obligations of a
future curved extension precisely.
\end{abstract}

\paragraph{Status and reading guide.}
The historical
\href{http://md.localhost/Users/matt/gcad/galacticad/docs/research/sfcc-algorithm-design.md}{SFCC design}
and the
\href{http://md.localhost/Users/matt/gcad/galacticad/docs/sfcc-meshing-algorithm.md}{implementation description}
provide project context, not premises of these proofs. In particular, the
protected-cell theorem changes the stopping and interface rules. It must not be
quoted as a proof that the original rules terminate, nor as a proof about the
Rust implementation. All results below use exact predicates and exact rational
constructions. There is no implementation change in this document.

\tableofcontents

\section{Target, inputs, and meanings of correctness}
\label{sec:target}

\subsection{Regularized solids, not every zero of a field}

For a subset $A\subseteq\R^3$, define
\[
  \reg(A)=\cl(\intr A).
\]
A solid is \emph{regular closed} when $A=\reg(A)$. The three operations in
this document are
\[
  A\mathbin{\cup^*}B=\reg(A\cup B),\qquad
  A\mathbin{\cap^*}B=\reg(A\cap B),\qquad
  A\mathbin{\setminus^*}B=\reg(A\setminus B).
\]
They are evaluated recursively in the supplied finite Boolean expression.
This is the regularized-solid convention described, for example, in the CGAL
Nef-polyhedra manual~\cite{cgal-nef}. Regularization removes isolated
lower-dimensional material; it does \emph{not} make every result manifold.

Our geometric target is
\[
  M=\partial S,
\]
where $S$ is that expression's regularized solid. A numerical field $f$ used
to represent the expression need not satisfy $f^{-1}(0)=M$.

Two elementary examples explain why this distinction is necessary.
If $d_A$ is negative inside $A$, the usual difference field for $A\setminus A$
is $\max(d_A,-d_A)=|d_A|$. It vanishes on $\partial A$, although
$A\setminus^* A=\varnothing$. Likewise, let two boxes share a full face and
occupy opposite sides of it. The minimum of their distance fields vanishes
on the shared face, while its relative interior lies inside the union solid.
Extracting all field zeros would retain a ghost sheet in either example.

\subsection{Standing assumptions}

Except in the orthogonal specialization, input primitives are bounded,
full-dimensional convex polyhedra given by finite rational linear
inequalities. Empty or lower-dimensional primitives can first be regularized
to the empty set. The expression uses only the three operations above.
There are no blends, offsets, transcendental surfaces, or unbounded root
complements. Consequently $S$ is bounded and regular closed.

We assume that nonempty $M$ is an embedded, closed topological $2$-manifold.
For a finite embedded polyhedral complex, this condition is decidable by the
edge-incidence and vertex-link checks described in Section~\ref{sec:features}.
Thus it can be an exact input-validation step, not an untestable promise.
The empty solid has the empty mesh as its output. Different connected
components, cavities, and arbitrary genus are allowed.

A root cube $Q=o+[0,L]^3$, with rational origin and positive rational side,
contains $S$ in its interior. In the generic-grid constructions, $L$ is a
positive integral power of two. A dyadic depth-$d$ cell has side
$h_d=L2^{-d}$. Cells are treated as \emph{closed} cubes; their interiors are
disjoint, and contacts on faces, edges, and vertices matter.
The notation $\Binf{p}{r}$ denotes the open ball for the maximum norm,
and $\Ninf{A}{r}=\{x:\dist_\infty(x,A)<r\}$ denotes its open metric
neighborhood about a set. A ``surface cell'' means a cube $C$ with
$M\cap C\ne\varnothing$, not merely a nonempty volume cube.

\subsection{What the output theorem asserts}

The output $T$ is a finite \emph{geometric simplicial complex}: its triangles
are nondegenerate, and two simplices intersect only in a common simplex or
not at all. Its realization, the union of all triangles, is $|T|$.
The exact theorems assert
\[
  |T|=M. \tag{1}\label{eq:exact}
\]
In particular they imply zero Hausdorff error, the correct topology, a
closed embedded $2$-manifold, and the identity ambient isotopy. Every true
corner is a mesh vertex, and every maximal straight crease is a union of
mesh edges with exactly its original geometric support. Artificial planar
and collinear subdivisions may remain in the mesh.

Equation~\eqref{eq:exact} is stronger than putting mesh vertices on the
surface. For a genuinely curved feature, finitely many straight mesh edges
cannot equal the feature as a set unless that feature is already piecewise
linear. A curved extension must specify approximation error and topology;
``exact feature vertices'' alone does not give exact curved features.

\section{Proof obligations that face sharing does not discharge}
\label{sec:obligations}

The following obligations are logically separate.
\begin{enumerate}[label=\textbf{O\arabic*.}]
\item \textbf{Semantics and completeness.} Recover the exposed boundary
  of the regularized solid and all of its genuine features. Candidate
  carrier intersections or finite sampled residuals do not prove this.
\item \textbf{Local geometry and topology.} Prove what $M\cap C$ is in
  every accepted cell, including components invisible at cube vertices.
\item \textbf{Termination.} Prove that the prescribed refinement and
  balancing rules reach accepted cells after finitely many steps.
\item \textbf{Complete interfaces.} Represent all of $M\cap\partial C$
  compatibly across adjacent cells, including pins and all subdivisions.
\item \textbf{Embedded filling.} Fill the specified patch without folds,
  crossings, omissions, or extraneous triangles, and respect its features.
\item \textbf{Assembly.} Show that the local triangulations fit as a
  geometric complex with the target surface and its orientation.
\end{enumerate}

Computing a face trace once is an effective way of enforcing identical
interface data; it does not prove the other obligations. Even ``every mesh
edge has two incident triangles'' allows two closed triangulated surfaces
to meet at a vertex, where the link has two cycles. Abstract manifoldness
also does not rule out intersecting geometric triangles. Finally, filling
an arbitrary spatial boundary loop by a fan can cross itself or leave the
intended surface. We will justify every fan by an exact plane or cone
description, rather than by its combinatorial form.

For approximate surfaces a useful gluing premise would be local
homeomorphisms that agree on their shared boundaries, together with a
separate embedding argument; an ambient isotopy needs still more control.
The polyhedral proofs avoid that gap by proving actual set equality in
every leaf and then checking simplex intersections.

\section{Exact S1: a finite complete boundary and feature compiler}
\label{sec:s1}

For S1 choose any auxiliary rational cube $Q_{\mathrm{aux}}$ strictly
enclosing the input primitives. It need not be the final octree root $Q$:
the latter is chosen after the genuine features have been compiled.
Changing the enclosing cube or its artificial clipping subdivisions
does not change the extracted solid boundary or its genuine features.

Let $\mathcal H$ be the finite set of distinct supporting planes of all
input inequalities, together with the auxiliary cube planes for clipping.
Write each plane as $a\cdot x=b$, with integer $a\ne0$ and integer $b$,
after clearing denominators. Coincident planes are identified.

\subsection{Chambers and occupancy}

The arrangement of $\mathcal H$ in $Q_{\mathrm{aux}}$ is a finite rational polytopal
complex. Its open $3$-dimensional chambers have constant signs for all
source inequalities. Each nonempty chamber has a rational interior witness.
Evaluating the Boolean expression at that witness labels the chamber
occupied or empty. On these witnesses ordinary Boolean evaluation and
recursive regularized evaluation agree, because no supporting plane is
encountered.

This construction is effective without an oracle. For example, enumerate
sign vectors, intersect the corresponding halfspaces with $Q_{\mathrm{aux}}$, and test
strict feasibility by maximizing a rational slack $\epsilon$ subject to
every chamber and auxiliary cube inequality having slack at least $\epsilon$.
A positive optimum certifies a full-dimensional chamber and yields a
rational witness. Arrangement faces and all incidences are obtained by
finite linear equalities and inequalities. This deliberately permits an
expensive finite enumeration; it is an existence and decidability argument.

Select each closed arrangement facet separating two chambers with
opposite occupancy labels. No auxiliary cube facet is selected, since
$S\subset\intr Q_{\mathrm{aux}}$.
Let $K$ be their union, with all their induced lower-dimensional faces.

\begin{theorem}[Complete exposed-boundary extraction]
\label{thm:arrangement-boundary}
The selected boundary complex satisfies $K=\partial S$.
\end{theorem}

\begin{proof}
Away from the union of source planes, primitive truth values are locally
constant. Induction on the Boolean expression shows the same for its
regularized evaluation. Regularization cannot change a truth value on an
open ball where that value is constant. Because the final set is regular
closed, it is precisely the closure of the union of its occupied chambers.
Indeed, the complement of finitely many planes is dense in every open
subset of $\intr S$.

In the relative interior of a selected facet there is an occupied chamber
on one side and an empty chamber on the other. Every neighborhood meets
both $S$ and its complement, so that facet interior is in $\partial S$;
closedness of the boundary includes the whole facet. Thus $K\subseteq M$.

Conversely, fix $x\in M$. A sufficiently small ball about $x$ meets no
arrangement plane that does not contain $x$. Its full-dimensional sectors
are therefore cones of the local arrangement. At least one sector is
occupied and at least one is empty: regular closedness supplies interior
solid points near $x$, while the boundary condition supplies exterior
points. The adjacency graph of these sectors across codimension-one faces
is connected. To see this, join interior points by a slightly perturbed
path in the ball that avoids the finite codimension-two arrangement
skeleton; it crosses only sector facets. A path from an occupied to an
empty sector must cross an unlike-label facet. The closure of this local
facet contains $x$, since the local arrangement is central at $x$.
It is part of a selected global facet. Hence $x\in K$.
\end{proof}

\subsection{Genuine strata, edges, and corners}
\label{sec:features}

Use the exact complex to remove artificial subdivisions. Define $G$ as
the set where the boundary fails to be locally an open subset of one
affine plane. It is a finite union of straight segments and their endpoints.
The closed planar strata $F_1,\ldots,F_N$ are the closures of the maximal
connected regular planar patches of $M\setminus G$. They are finite unions
of rational polygons, on carrier planes $P_1,\ldots,P_N$, and satisfy
\[
  M=\bigcup_i F_i,\qquad
  F_i\cap F_j\subseteq G\ (i\ne j),\qquad
  \partial_{P_i}F_i\subseteq G. \tag{2}\label{eq:strata}
\]
Here $\partial_{P_i}$ is the boundary within the carrier plane. Coplanar
facets are merged across artificial edges where the surface has a planar
germ. Merely touching at a true corner is not a reason to identify two
different regular patches.

Subdivide $G$ at every change in direction or incident face germs, then
merge collinear degree-two subdivisions with unchanged flanking germs.
Write $E$ for the resulting maximal closed straight edges and $V$ for
their endpoints, including all true geometric corners. At every point
of $\relint e$, for $e\in E$, exactly two noncoplanar planar half-sheets
meet in a fixed wedge. The two incident closed strata are denoted
$F_{e,1}$ and $F_{e,2}$. An edge endpoint cannot be removed merely because
two sampled normals are close: removal requires exact equality of the
local geometry.

All these decisions are finite arrangement and incidence decisions.
Locally, inspect the cyclic occupied/empty sectors and the selected planar
facets. This also checks the manifold premise: first triangulate the
embedded boundary arrangement conformingly; every mesh edge must have
exactly two incident triangles, and every vertex link must be a single
cycle. Since the arrangement already supplies a geometric complex,
there is no additional geometric self-intersection assumption hidden in
those combinatorial tests. Regularized Boolean solids can fail them;
for example, two boxes joined only along an edge are outside the theorem.

\begin{remark}[What exact S1 replaces]
This S1 uses exposed arrangement facets, not all pairwise carrier
intersections. It consequently includes Boolean-created seams and corners,
removes hidden carrier intersections and ghost sheets, and merges
geometrically artificial subdivisions. It is an exact polyhedral boundary
compiler. No claim is made that a sampled feature tracer implements it.
\end{remark}

\section{Result A: orthogonal-box completion with simple cells}
\label{sec:orthogonal}

Suppose now that every input primitive is an axis-aligned rational box.
Let $X,Y,Z$ be the finite sets of its support coordinates. These partition
the root into rectangular chambers. Rational interior witnesses and exact
Boolean occupancy select the exposed rectangles just as in
Theorem~\ref{thm:arrangement-boundary}. Their closures are the whole of $M$.

\begin{theorem}[Orthogonal simple-cell theorem]
\label{thm:orthogonal}
For a regularized Boolean combination of bounded rational axis-aligned
boxes with embedded closed manifold boundary, there is an effective
rational dyadic root grid and a finite, optionally $2{:}1$ balanced octree
such that every leaf with nonempty surface intersection contains exactly
a smooth disk, an edge disk, or a one-corner disk. Complete shared face
arcs and the exact fillings below
produce a geometric simplicial complex satisfying~\eqref{eq:exact}.
The sufficient stopping certificate is: \emph{at most one support plane
per coordinate axis intersects the cell}. This is a new exact certificate,
not a consequence of corner signs or sampled normal tests.
\end{theorem}

\subsection{An all-depth grid and a finite stopping depth}

Choose an odd prime $q$ dividing none of the denominators of the input
coordinates in lowest terms. Choose integers $b_i$ below the geometry,
put $o_i=b_i+1/q$, and take a sufficiently large power-of-two side $L$.
For any support coordinate $a=p/s$ with $q\nmid s$,
\[
  \frac{a-o_i}{L}
  =\frac{pq-b_i sq-s}{sqL}.
\]
Its numerator is $-s$ modulo $q$, hence its reduced denominator still
contains the odd factor $q$. It cannot be dyadic. Thus no support plane
coincides with any grid plane at any depth.

Let $\delta$ be the minimum positive gap between two distinct support
coordinates on the same axis, omitting axes with fewer than two. Choose
finite $D$ with
\[
  L2^{-D}<\delta. \tag{3}\label{eq:orth-depth}
\]
If there are no such gaps, any depth with the stated one-plane certificate
suffices. At depth $D$ every closed cell meets at most one support plane
per axis. This certificate is inherited by children. Refine cells until
it holds and optionally add face or edge $2{:}1$ balancing. Balancing
cannot create depth $D+1$: a split forced solely by imbalance is applied
to a coarser neighbor of a finer existing cell, and does not exceed that
cell's depth. Induction bounds every created depth by $D$.

\subsection{Every accepted cell is a truncated cone}

Fix an accepted cell $C$. Define $v\in\intr C$ by using the unique support
coordinate on each axis that has a cut and the cell midpoint on each
missing axis. All primitive signs in $C$ depend only on the signs of
$x_i-v_i$ for the present cuts. Missing axes do not change any label.
Therefore the occupied set, and hence its boundary, are invariant along
rays from $v$ until those rays exit $C$. There is an exact polyhedral
cone $\mathcal K_v$ such that
\[
  M\cap C=\mathcal K_v\cap C.
\]
If this set is nonempty, $v\in M$. The link of the boundary cone on a
small sphere around $v$ is a single simple cycle: at each link point the
manifold premise gives degree two, and multiple link cycles would make
the cone tip a nonmanifold point. Radial projection from $v$ to the cube
boundary is a homeomorphism on directions. Thus $M\cap C$ is a disk,
with boundary $M\cap\partial C$.

\subsection{The link gives the three cell classes without a lookup table}

Extend the sign description to eight octants by duplicating labels on
missing axes. The coordinate great circles divide a sphere into eight
triangles. Their one-skeleton is the octahedron graph, with vertices
$\pm e_1,\pm e_2,\pm e_3$. The link consists of edges separating octants
with opposite labels.

At any octahedron vertex, four octants occur cyclically. Alternating
labels would produce four incident link edges, corresponding to a
nonmanifold coordinate ray. All used vertices therefore have degree
two, and the manifold condition at $v$ leaves one cycle. Along a link
edge the surface lies in one coordinate plane. A \emph{turn} at a link
vertex means that the two incident link edges use different coordinate
great circles; exactly these turns represent genuine sharp rays.

If there are no turns, the cycle is one coordinate great circle and
the patch is planar. If all turns are on one coordinate axis, they are
the two antipodal vertices of that axis: a change of great circle can
return only at their other intersection. The patch has one continuing
straight crease. In all other cases sharp rays use at least two axes,
and $v$ is a genuine corner incident to all of them. This proves the
three classes directly; it assumes no maximum valence beyond what the
orthogonal geometry itself permits.

For an independent finite cross-check, enumeration of the $2^8$ occupancy
masks gives the following counts. This table is not a premise of the proof.
\begin{center}
\begin{tabular}{lr}
\toprule
Local configuration & Number of masks\\
\midrule
No boundary & 2\\
Plane & 6\\
Continuing edge & 24\\
Corner of degree $3$ & 64\\
Corner of degree $4$ & 24\\
Corner of degree $6$ & 8\\
At least one nonmanifold coordinate ray & 120\\
Degree-two rays but disconnected vertex link & 8\\
\midrule
Total & 256\\
\bottomrule
\end{tabular}
\end{center}

\subsection{Faces, witnesses, and exact filling}

Each cell face has at most two perpendicular coordinate cuts. Its sectors
have a cube-corner witness each. With two cuts, the only ambiguous
checkerboard labeling would give a nonmanifold ray perpendicular to the
face; it is excluded by the manifold premise. The face boundary trace
is consequently empty, a straight arc, or one L-shaped arc. With fewer
cuts the same conclusion is immediate. Every three-dimensional sign
sector has a cube-corner witness as well. There are no hidden sheets or
closed face loops under this particular support-coordinate certificate.

Compute each shared face arc once, including crease pins and all
subdivisions required by neighboring leaves. A full-depth-$D$ interface
subdivision is a finite, simple conforming choice, although much less
subdivision normally suffices.

For a corner cell, cone each consecutive boundary segment $[a,b]$ to
$v$. The segment lies on a single actual support plane through $v$,
and every radial segment $[v,p]$, $p\in[a,b]$, belongs to the exact
surface. Thus the triangle $[v,a,b]$ lies in $M$. Each ray from an
interior point of a convex cube has a unique exit, so these triangles
cover the patch and have disjoint relative interiors. Smooth and edge
cells can use the same argument or planar fans on their one or two
convex polygonal pieces. Shared subdivisions are retained.

These local triangulations agree across faces. The precise assembly
argument in Section~\ref{sec:assembly} applies verbatim and proves
Theorem~\ref{thm:orthogonal}.

\section{Result B: fixed-grid balanced refinement need not terminate}
\label{sec:obstruction}

\subsection{The rule being refuted}

Consider the following rules on a fixed dyadic root:
\begin{enumerate}[label=(\roman*)]
\item Each accepted leaf meets at most one genuine feature edge, or
  contains exactly one corner and only feature edges incident to it.
\item Leaves sharing positive face area differ in depth by at most one
  (face $2{:}1$ balance).
\item Refinement alone is used to make every leaf accepted.
\end{enumerate}
Feature contacts below are transverse, so interpreting ``meets'' via
interior crossings instead of closed-cube contacts does not remove the
obstruction. The claim refuted is universal finite acceptance on the
specified grid, not the possibility of meshing a tetrahedron by other rules.

\begin{theorem}[A rational tetrahedron obstruction]
\label{thm:tetrahedron}
On the root $[0,1]^3$, let
\[
  v=(1/3,1/3,1/3),
\]
and define
\begin{align*}
 b_1&=v+(0,1,-61)/366,\\
 b_2&=v+(1,0,-61)/366,\\
 b_3&=v+(-3,-2,-61)/366.
\end{align*}
The boundary of $\conv\{v,b_1,b_2,b_3\}$ is an embedded closed rational
polyhedral $2$-manifold, but no finite octree can satisfy all three rules
above for this boundary.
\end{theorem}

\begin{proof}
All three base vertices have $z=1/6$. In the base plane, the determinant
of $(1,-1)$ and $(-3,-3)$ is $-6$, so the base is nondegenerate and the
apex is off its plane. All four vertices are strictly inside the root.

The apex leaf cannot have depth $0$ or $1$, since at either depth its
cell contains all four corners. Suppose its depth is $d\ge2$ and let
$h=2^{-d}$. Each coordinate of the apex has fractional position
\[
  \{2^d/3\}\in\{1/3,2/3\}
\]
within that cell. Its distance from every side face is at least $h/3$,
and its projected distance from either midline of a horizontal face is
$h/6$.

Write $z_0=h\lfloor 2^d/3\rfloor$ for the lower face height and
$\Delta=1/3-z_0\in\{h/3,2h/3\}$. Along each of the three apex edges,
horizontal slope per unit downward displacement is one of
\[
  (0,1)/61,\qquad(1,0)/61,\qquad(-3,-2)/61.
\]
The horizontal drift before reaching $z=z_0$ is at most
\[
  \frac{3\Delta}{61}\le\frac{2h}{61}<\frac h6.
  \tag{4}\label{eq:tetra-drift}
\]
Thus all three edges exit the lower face, none can reach a side first,
and all three exit points lie strictly inside the same quarter of the
lower face. Since $z_0\ge1/4>1/6$, these crossings precede the base
vertices and are interior points of their edges.

Face $2{:}1$ balance permits only depth-$d-1$, depth-$d$, or depth-$d+1$
neighbors across that face. Dyadic alignment implies that its chosen
open quarter is contained in one neighboring leaf $N$; no permitted
neighbor boundary can split that quarter. All three edges therefore
enter the interior of $N$. The neighbor cannot contain $v$, being below
$z_0$. No other tetrahedron vertex is incident to all three apex edges.
It follows that $N$ has neither one feature edge nor one corner with
only its incident edges. This contradicts acceptance.
\end{proof}

\subsection{This example is not a grid-coincidence pathology}

All four vertices have nondyadic coordinates. Each apex direction
\[
  w\in\{(0,1,-61),(1,0,-61),(-3,-2,-61)\}
\]
has three components pairwise distinct modulo $3$.
If two coordinates of a point $p=v+tw$
were dyadic, then
\[
  w_jp_i-w_ip_j=(w_j-w_i)/3
\]
would be both dyadic and a rational with reduced denominator divisible
by $3$, an impossibility. Hence no apex edge meets a dyadic grid edge.

The three base edges are in $z=1/6$. Their projected lines respectively
satisfy
\[
  x+y=245/366,\qquad x-y=-1/366,\qquad x-2y=-121/366.
\]
These right-hand sides are nondyadic, so their projections contain no
point with two dyadic coordinates. The base edges avoid every grid edge
as well. A feature line parallel to a grid face does not meet that face;
all actual feature/face crossings are transverse. The obstruction comes
from clustered directions and balancing, not exact incidences or
tangencies.

\subsection{Two further obstructions and their scope}

\begin{proposition}[The $24$-neighbor obstruction]
\label{prop:valence}
Under face $2{:}1$ balance and generic feature/face intersections, a
one-corner cell with at least $25$ distinct incident straight edges
cannot have every face neighbor satisfy the stated simple-cell rule.
\end{proposition}

\begin{proof}
An accepted corner cell contains no other edge endpoint. Each incident
edge exits through a face interior. There are at most four neighbors
per face, hence at most $24$ face-neighbor leaves. Two edges enter one
neighbor by the pigeonhole principle. They have no common corner other
than the original corner: distinct straight segments sharing both
endpoints would be the same geometric edge. The neighbor does not
contain the original corner and therefore cannot be accepted.
\end{proof}

Rational pyramids over rational convex $25$-gons provide such manifold
corners in generic position. The tetrahedron theorem shows that degree
at most $24$ is far from sufficient; degree three already fails.

For the tetrahedron apex cell, the exact trace on its lower face is a
small closed triangle entirely in that face's interior. All cube
corners, and all face corners, are outside the solid. A literal
Marching-Cubes complex-face test that rejects sign-change-free face
contours consequently rejects this cap at every depth. A feature-aware
cone certificate can accept it, but that is a different certificate.
The topology result of Varadhan et al.~\cite{varadhan} does not supply
termination for this modified setting merely by being cited.

Finally, a fixed positive threshold on the angle between a straight
feature and a grid face cannot be repaired by shrinking cells: the
line direction and face normal remain unchanged. The same observation
applies to any fixed angular exclusion for an unchanged carrier plane.
Exact nonzero transversality is weaker than a uniform positive angle
cutoff. None of the positive results below requires such a cutoff.

\section{Result C: unrestricted exact arrangement completion}
\label{sec:fallback}

\begin{theorem}[Finite arrangement fallback]
\label{thm:fallback}
For the rational-polyhedral input class, any finite octree admits a
finite exact arrangement completion whose triangle realization is $M$.
This includes arbitrary local feature arrangements and surface pieces
that coincide with leaf faces. The construction terminates without
the simple-cell stopping rule.
\end{theorem}

\begin{proof}
Take all source support planes and all supporting planes of all leaf
faces. There are finitely many. Their global arrangement in $Q$ is a
finite rational polytopal complex. Label its $3$-chambers and select
unlike-label facets. Adding leaf planes changes only subdivisions, so
Theorem~\ref{thm:arrangement-boundary} identifies their union with $M$.

The arrangement ensures that adjacent selected convex polygons have
exactly the same subdivided common edges. Insert one canonical midpoint
on each arrangement edge. For each polygon choose a rational interior
point, for example the mean of its extreme vertices, and connect it to
each consecutive boundary subsegment. These are the usual facewise
barycentric fans. Convexity proves that they tile the polygon with
nonoverlapping nondegenerate triangles. Shared edge subdivisions agree,
so the result is a geometric simplicial complex with realization $M$.

If a selected facet lies in a grid face, it is a two-dimensional
interface patch, not a contour arc. Emit it once with a deterministic
ownership convention. All incidences are still inherited from the one
global arrangement, so this coincidence causes no geometric ambiguity.
Every enumeration, feasibility decision, and triangulation is finite.
\end{proof}

One may stop octree refinement at any fixed depth and invoke this
fallback. It is an exact boundary-representation mesher organized
around an octree, and it permits many patches or features in one leaf.
It is \emph{not} a proof of termination for the original smooth/edge/corner
rule. The next theorem keeps those three volume-cell templates by
changing protection and interface handling instead.

\section{Result D: protected-neighborhood completion}
\label{sec:protected}

\begin{theorem}[Protected polyhedral SFCC]
\label{thm:main}
Let $S$ satisfy the rational-polyhedral and embedded-manifold assumptions
of Section~\ref{sec:target}. There is an effective exact-arithmetic
algorithm that returns a finite octree with frozen corner and edge
leaves, complete compatible face graphs, and a finite rational
geometric triangulation $T$ such that $|T|=\partial S$.
Every leaf with nonempty surface intersection has one of these surface
patches:
\begin{enumerate}[label=(\alph*)]
\item a disk cut from one affine plane;
\item a disk formed by two convex planar polygons meeting along one
  straight crease segment;
\item an exact conical disk about one true corner, containing only its
  incident feature rays, with unrestricted corner valence.
\end{enumerate}
All true corners and straight creases are preserved exactly. The
algorithm terminates at three finite, effectively chosen depths
$D_0\le D_1\le D_2$. It does not impose $2{:}1$ balancing across
protected volume boundaries, and it does not subdivide a protected
volume leaf to simplify its boundary graph.
\end{theorem}

\subsection{A rational generic grid valid at every depth}
\label{sec:generic-grid}

We need the following incidence exclusions:
\begin{enumerate}[label=(\roman*)]
\item no feature vertex lies on a grid face plane;
\item no feature line meets a grid edge line;
\item no surface carrier plane contains a grid vertex.
\end{enumerate}
These also exclude a carrier plane coincident with a grid face and
a feature segment lying in a grid face. They do not demand a lower
bound on nonzero intersection angles.

\begin{lemma}[Effective all-depth rational genericity]
\label{lem:grid}
A rational origin and power-of-two side $L$ satisfying these exclusions
at all dyadic depths can be constructed by a finite calculation.
\end{lemma}

\begin{proof}
Build a finite list of rational affine forbidden planes containing:
all primitive support planes; all three coordinate planes through every
feature vertex; and the following projected-line planes. For a feature
line $x=p+tu$, and each coordinate pair $i,j$ with $(u_i,u_j)\ne(0,0)$,
include
\[
  u_jx_i-u_ix_j=u_jp_i-u_ip_j. \tag{5}\label{eq:projected-plane}
\]
If both components are zero, coordinate planes through an endpoint
already exclude the corresponding constant projection. Here ``source
support planes'' means primitive planes, not auxiliary clipping planes
or final octree root planes. The latter are excluded from this list:
root planes intentionally contain grid vertices. Each listed plane has
an equation $a\cdot x=c$ with all coefficients integral, including the
right-hand side $c$, and with nonzero integer normal $a$, after clearing
denominators. Let the number of planes be $m$.

Choose an odd prime $q>m$ such that every listed normal remains nonzero
modulo $q$. Only finitely many primes are excluded, so trial search
terminates. For each normal, the set of $r\in\Fq^3$ satisfying
$a\cdot r=0$ has $q^2$ elements. The union of the $m$ forbidden sets
has at most $mq^2<q^3$ elements. Exhaustive finite search therefore
finds $r$ with $a\cdot r\ne0$ for every listed normal.

Represent $r$ by integers in $\{0,\ldots,q-1\}$ and put
$o=b+r/q$, where $b\in\Z^3$ lies sufficiently far below the geometry.
Choose an integral power of two $L$ with $S\subset\intr Q$. If a grid
vertex $x=o+Lk/2^d$ were on a forbidden plane, multiplying its equation
by $q2^d$ and reducing modulo $q$ would give
\[
  2^d a\cdot r=0\pmod q,
\]
contradicting the choice of $r$. This proves avoidance at every depth
simultaneously, including vertices of the infinite extension of the grid.

A feature vertex on a coordinate grid plane would put a grid vertex
on its forbidden coordinate plane. A feature line meeting a grid edge
would have two coordinates equal to grid coordinates; substituting
these into~\eqref{eq:projected-plane} would put a grid vertex on that
forbidden plane. The constant-projection case is covered by the
endpoint coordinate planes. Carrier-plane/grid-vertex incidences are
excluded directly. This proves all three claims.
\end{proof}

\subsection{Exact conical neighborhoods of corners}
\label{sec:corners}

For $v\in V$, each source plane $a\cdot x=b$ not passing through $v$
keeps its sign in a sufficiently small cube. Specifically choose
rational $r_v>0$ so that
\[
  r_v<\min_{a\cdot v\ne b}
    \frac{|a\cdot v-b|}{2\norm{a}_1}. \tag{6}\label{eq:corner-radius}
\]
Omit the minimum when its index set is empty. Also take $r_v$ smaller
than suitable positive distances to other vertices, nonincident
closed features, and $\partial Q$. These distances are rationally
computable as shown in Section~\ref{sec:effective}.

\begin{lemma}[Conical corner germ]
\label{lem:corner-cone}
There is an exact boundary cone $\mathcal K_v$ centered at $v$ such that
$M\cap\Binf{v}{r_v}=\mathcal K_v\cap\Binf{v}{r_v}$. Its link is a
single embedded polygonal cycle. If a closed cube $C$ contains $v$ in
its interior and lies in this neighborhood, then $M\cap C$ is a radial
disk with boundary $M\cap\partial C$.
\end{lemma}

\begin{proof}
For a nonthrough plane, the change of $a\cdot x-b$ from its value at
$v$ is bounded by $\norm a_1\norm{x-v}_\infty$, hence its sign is
fixed by~\eqref{eq:corner-radius}. Equations of through planes are
homogeneous about $v$. Every chamber label, and therefore the
regularized boundary, is constant along each ray within this
neighborhood. This gives the exact cone.

Intersect the cone with a small sphere about $v$. It is a finite
embedded graph of arcs of great circles. Away from its vertices a
surface sheet gives an arc; at every graph vertex the surface-manifold
condition forces degree two. The link is thus a union of circles.
If it had more than one circle, a punctured sufficiently small
neighborhood of $v$ in $M$ would have more than one component, whereas
a punctured disk is connected. Hence it is one circle.

For a convex cube containing $v$ in its interior, every ray has a
unique boundary exit, continuously dependent on direction. Radial
projection sends the link homeomorphically to $M\cap\partial C$;
adjoining all radial segments makes a cone over a circle, hence a disk.
Convexity of the solid or of the spherical region enclosed by the link
is not required. Reflex and high-valence corners are included.
\end{proof}

By Lemma~\ref{lem:grid}, every feature vertex belongs to the interior
of a unique cell at each depth. Choose a common finite $D_0$ so that
the owning cubes $C_v$ lie in their conical neighborhoods and have
pairwise disjoint closures. Such a depth exists because cell diameter
tends to zero and distinct vertices have positive separation.
\emph{Freeze these volume cells.} Every later step retains them as
leaves, even when adjacent leaves become much smaller.

\subsection{Exact wedge neighborhoods of clipped edges}
\label{sec:edges}

For each feature edge remove the interiors of its endpoint corner cubes:
\[
  e^*=e\setminus\bigcup_{v\in V}\intr C_v.
\]
The corner choices ensure that this is a closed segment separated from
the true endpoints, and that no nonincident corner cube cuts it.
Different such clipped edges are disjoint compact sets. Empty clipped
pieces, if any redundant case is present, can be omitted.

Let $W_e$ be the complete straight wedge obtained by extending the two
local incident half-planes along the carrier line of $e$. Define
\[
  B_e=V\ \cup\!\bigcup_{f\in E,\,f\ne e}f
    \ \cup\!\bigcup_{F_i\notin\{F_{e,1},F_{e,2}\}}F_i.
  \tag{7}\label{eq:edge-bad}
\]
All sets here are closed and bounded. By the exact stratification,
$e^*\cap B_e=\varnothing$, so
\[
  d_e=\dist_\infty(e^*,B_e)>0.
\]
One may include $\partial Q$ in $B_e$ to keep every chosen neighborhood
within the root. If a bad set is empty, choose any positive rational
bound in its place. Take $0<\rho_e<d_e/4$ and put
\[
  U_e=\Ninf{e^*}{\rho_e}
      =e^*+(-\rho_e,\rho_e)^3.
\]
It is a convex open tube.

\begin{lemma}[Uniform exact edge wedge]
\label{lem:edge-wedge}
The exact boundary in this tube is $M\cap U_e=W_e\cap U_e$.
\end{lemma}

\begin{proof}
The tube meets no vertex, other feature edge, or nonincident stratum.
Its intersection with the carrier line is a connected interval within
the interior of the same feature edge: moving beyond that edge would
require reaching an endpoint in $V$.

In either incident carrier plane, removing that line divides the tube
into two convex, hence connected, relative-open half-plane pieces.
By~\eqref{eq:strata}, membership in the incident closed stratum cannot
change within one such piece without meeting a planar boundary in
$G$. No such boundary is present. Membership is therefore constant
on each piece. The known two-face wedge germ at any point of $e^*$
determines exactly which half is present on each plane. No other
stratum is present, so the union is exactly $W_e\cap U_e$.
\end{proof}

Choose common $D_1\ge D_0$ so that
\[
  h_{D_1}<\min_e\rho_e. \tag{8}\label{eq:edge-depth}
\]
Refine all noncorner cells to depth $D_1$, without refining any $C_v$.
\emph{Freeze every remaining cell whose closed cube meets $G$.}
Every such cell $C$ contains a point $p\in e^*$, so
$C\subset\Ninf{e^*}{\rho_e}$ by~\eqref{eq:edge-depth}.
It meets no other edge or corner.

Genericity excludes a feature line touching only a cube edge or vertex
or lying in a cube face. Thus $e\cap C$ is a nontrivial segment through
the cube interior. Each half-plane of $W_e$ cuts $C$ in a convex planar
polygon with positive relative area. They meet exactly in that crease
segment, on both polygon boundaries. Their union is an exact disk.
This is the edge-cell template, even if the two planes meet at a very
small or reflex dihedral angle.

\begin{remark}[An alternative exact clearance certificate]
One can instead form the closed rational polyhedral bad set
\[
  D_e=\cl\bigl((M\symdiff W_e)\cap Q\bigr).
\]
Every point of $e^*$ has a wedge neighborhood, so $e^*\cap D_e$ is
empty and its distance is positive. This is an effective exact set
comparison using the S1 arrangement and wedge planes, including
lower-dimensional strata. It is not enough to compare only sampled
points or only positive-area pieces. The explicit bad set
\eqref{eq:edge-bad} gives a direct geometric proof and usually a more
transparent certificate.
\end{remark}

\subsection{Why closed contacts protect an open neighborhood}
\label{sec:open-neighborhood}

\begin{lemma}[Open neighborhood from closed-contact freezing]
\label{lem:open-neighborhood}
After the corner and edge steps, the union of frozen cells contains an
ambient open neighborhood of $G$.
\end{lemma}

\begin{proof}
Let $R$ be the union of the \emph{closed} cells that remain unfrozen.
It is a finite union of compact cubes. Every remaining cube has empty
intersection with $G$, since all closed-contact cubes were frozen.
Consequently $R\cap G=\varnothing$. If $R$ is nonempty, compactness
gives $\dist_\infty(R,G)>0$; if it is empty, no separation is needed.
The graph also has positive distance from $\partial Q$. A sufficiently
small ambient neighborhood of $G$ therefore stays in $Q$ and avoids
$R$. Since all cells together cover $Q$, that neighborhood lies in the
union of frozen cells.
\end{proof}

The word ``closed'' is essential. If only cells with an interior
feature intersection were frozen, a cell touching a feature on its
boundary could remain. Then $R\cap G$ need not be empty, and the
positive-distance argument for the smooth remainder would fail.

\subsection{Positive separation makes the remainder planar}
\label{sec:smooth}

For each closed stratum let
\[
  A_i=F_i\cap R,\qquad
  B_i=G\cup\bigcup_{j\ne i}F_j.
\]
Omit empty $A_i$. Equation~\eqref{eq:strata} and
Lemma~\ref{lem:open-neighborhood} imply $A_i\cap B_i=\varnothing$.
Both sets are compact, so choose positive rational
\[
  s_i<\tfrac14\dist_\infty(A_i,B_i). \tag{9}\label{eq:smooth-clearance}
\]
As before, including $\partial Q$ or assigning a bound to an empty
bad set is harmless.

\begin{lemma}[Planar remainder]
\label{lem:smooth}
If a remaining descendant cube $C\subseteq R$ has side less than
$s_i$ and meets $F_i$, then $M\cap C=P_i\cap C$. In particular every
nonempty sufficiently fine remaining cell contains just one convex
planar disk.
\end{lemma}

\begin{proof}
Choose $p\in C\cap F_i\subseteq A_i$. The cube lies inside
$\Binf{p}{s_i}$. That ball contains no $G$ and no other closed stratum.
Its intersection with $P_i$ is convex and connected. The subset
belonging to $F_i$ is nonempty, closed, and relatively open there,
because its planar boundary is contained in $G$. It is therefore the
entire planar section. No other surface is present. Thus
$M\cap C=P_i\cap C$.

The generic grid excludes a plane touching the cube only at a vertex
or along an edge: either contact would contain a grid vertex. It
also excludes a plane coincident with a cube face. A nonempty
intersection is consequently a convex polygon with positive area.
\end{proof}

Refine the remaining cells to common $D_2\ge D_1$ with
$h_{D_2}<\min_i s_i$. Protected cells remain unchanged. This proves
finite geometric simplification of the remainder. It is not a corner-sign
argument: a thin slab can have two parallel sheets through a cube and
all eight corners outside, with no crease in the cube. The positive
distance from other closed strata is what excludes that example.

An equivalent exact certificate uses
$D_i=\cl((M\symdiff P_i)\cap Q)$. Each $A_i$ is disjoint from $D_i$,
and a positive distance to $D_i$ certifies the same local plane equality.
Again this comparison must respect lower-dimensional strata.

\subsection{Every bound and decision is effective}
\label{sec:effective}

For two nonempty disjoint compact rational polytopes $A,B$, their
$L^\infty$ distance is the optimum of the rational linear program
\[
\begin{aligned}
  \text{minimize }&t\\
  \text{subject to }&x\in A,\quad y\in B,\quad t\ge0,\\
  &-t\le x_k-y_k\le t\quad(k=1,2,3).
\end{aligned} \tag{10}\label{eq:lp-distance}
\]
Equalities handle segments and planar polygons without artificially
thickening them. The optimum is rational and positive when the compact
sets are disjoint. For finite unions of rational polytopes take the
minimum over nonempty component pairs. All sets used above have that
form, including clipped edges, $R$, and closed strata.

The construction of $r_v$, separation of corner cells, and choices of
$D_0,D_1,D_2$ are thus finite searches with rational comparisons. Empty
index sets impose no restriction. Distances can be extremely small,
and the resulting depths large, but positivity is proved rather than
assumed. Neither a global angle bound nor a floating-point tolerance
is used to infer separation.

\section{Interfaces, local triangulations, and the final assembly}
\label{sec:assembly}

\subsection{Complete shared face graphs without volume balancing}

The finite leaf octree just constructed can have arbitrarily different
depths on opposite sides of a protected cell face. The interface must
accommodate this without subdividing that protected volume.

One explicit finite solution is to subdivide \emph{every actual leaf
face} into depth-$D_2$ grid squares. This is a subdivision of the
two-dimensional interfaces only; it adds no children to corner or edge
volume cells. Square identities are canonical by axis, grid coordinate,
and in-plane indices. On each square $F$, compute
\[
  \Gamma_F=M\cap F
\]
once, by exact clipping of the exposed S1 boundary facets. Since a
carrier plane never coincides with a grid face, this is a finite
straight-segment graph, with no two-dimensional patches.

Node every segment at all its intersections, feature pins, and shared
interface subdivisions. Merge coincident segments and identify points
by exact rational position, retaining feature incidence labels.
Subdivisions on the common edges of interface squares are reconciled
globally. Both incident volume cells then consume the same graph;
a coarse face consumes the union of its constituent square graphs.

The representation must admit multiple arcs, closed face loops, and
multiple roots on a coarse grid edge. It cannot rely on a perimeter
sign change to seed every component. The tetrahedron's corner cap
gives a complete loop on a single coarse face. Its correct corner
fill is a cone from the apex; triangulating the loop \emph{inside the
face plane} would instead insert a spurious lid.

Uniform depth-$D_2$ interface subdivision is only an existence
construction. A common adaptive face refinement works too if it
gives exactly the same complete graph and all shared subdivisions.
Fewer interface squares change the cost, not the proof obligations.

\subsection{Three exact filling templates}

\paragraph{Corner leaf $C_v$.}
Its complete boundary graph is one simple polygonal loop by
Lemma~\ref{lem:corner-cone}. Split it at every change of carrier plane,
crease pin, cell-face transition, and required interface node. For
each nonzero consecutive segment $[a,b]$, emit $[v,a,b]$.
Every such segment is in a support plane through $v$; radial cone
equality puts the whole triangle on $M$. Because $v\in\intr C_v$
and $[a,b]$ lies on a cube face, the triangle is nondegenerate. Unique
radial exits prove exact coverage and nonoverlap, including for
reflex or nonconvex boundary loops. Crease rays occur as mesh edges
from the exact corner to the appropriate boundary pins.

\paragraph{Edge leaf $C$.}
The exact patch is the union of the two convex planar polygons of
Lemma~\ref{lem:edge-wedge}, joined along $e\cap C$. Add every prescribed
boundary node from the interface graph and a common subdivision of
the crease if one is needed. In each polygon choose a rational point
in its relative interior and cone it to the consecutive boundary
subsegments. The mean of the polygon's extreme vertices is suitable.
Both polygons have interior points in $\intr C$; the fan points may
and do lie there. Convexity proves exact coverage and nondegeneracy,
and their only common geometry is the prescribed crease chain.
Equivalently, an appropriately subdivided fan from a point inside
the crease can be used; a separate interior point for each planar
piece avoids needing that variant in the proof.

\paragraph{Smooth leaf $C$.}
The exact patch $P_i\cap C$ is one convex polygon. Retain all of its
boundary nodes, choose a rational relative-interior point in
$\intr C$, and fan to the boundary subsegments. Empty cells emit
nothing. No quality improvement, vertex perturbation, or face-node
deletion is part of these exact templates.

\begin{lemma}[Exact compatible assembly]
\label{lem:assembly}
The three templates with the complete shared graphs form a finite
rational geometric simplicial complex $T$ satisfying $|T|=M$.
Its facets can be oriented consistently from the occupied side of
the boundary to the empty side.
\end{lemma}

\begin{proof}
Each template gives a nondegenerate, nonoverlapping triangulation of
exactly $M\cap C$, using the prescribed subdivision of
$M\cap\partial C$. For a planar fan this follows from convexity;
for a corner fan it follows from the unique radial exit. In an edge
cell the two planar fans meet only along their identical crease chain.

Within a leaf, simplices therefore meet face-to-face. A triangle's
intersection with the leaf boundary is its prescribed boundary base
edge or its endpoints: the fan point is in the volume interior.
Two distinct leaves have disjoint interiors. All possible
inter-leaf triangle intersections lie on interfaces, where the
same noded segment graph is used. Two such triangle bases are
identical, disjoint, or meet at one common node. Therefore triangles
from different leaves also intersect only in common simplices.
The globally identified vertices and all simplices constitute a
geometric simplicial complex.

The leaves cover $Q$, so
\[
  |T|=\bigcup_C(M\cap C)=M.
\]
In particular all edge and vertex links are those of the assumed
embedded manifold, not just a complex passing an edge-count check.
At a regular facet the chamber occupancy determines the outward
orientation. These choices are compatible across the boundary of
the regular solid, so adjacent triangles induce opposite directions
on each common edge. This includes interfaces and crease edges.

Corner fan centers are the exact true corners. In all three
templates, crease portions are retained as mesh edges and shared
pins are identified. Their union is exactly $G$, with its original
corner incidence and geometry.
\end{proof}

\subsection{A finite algorithm and its dependency chain}

\begin{algorithm}[Protected polyhedral SFCC]
\label{alg:protected}
Given the rational hard-CSG input:
\begin{enumerate}
\item Construct the support arrangement; label chambers; extract
  exposed boundary facets. Check manifold incidence and links.
  If empty, return the empty mesh; if nonmanifold, report that the
  theorem's input condition fails.
\item Merge artificial planar and collinear subdivisions to obtain
  $F_i$, $G$, $E$, and $V$ with exact incidence and side labels.
\item Construct the all-depth generic rational grid of
  Lemma~\ref{lem:grid}.
\item Compute corner radii. Choose common $D_0$ with disjoint owning
  corner cubes contained in their cone neighborhoods; freeze them.
\item Compute clipped edges and their wedge clearances. Refine only
  noncorner cells to $D_1$ satisfying~\eqref{eq:edge-depth}; freeze
  every cell whose \emph{closed} cube meets $G$.
\item Form the finite closed remainder $R$. Compute the planar
  clearances~\eqref{eq:smooth-clearance}; refine only remaining cells
  to $D_2$. The protected volume leaves remain frozen.
\item Build a finite common interface subdivision. Compute and node
  each complete exact shared face trace once.
\item Apply the corner, edge, and smooth filling templates; identify
  common vertices and edges; orient by occupied-to-empty sides.
\end{enumerate}
\end{algorithm}

\begin{proof}[Proof of Theorem~\ref{thm:main}]
Exact S1 and its finite checks follow from
Theorem~\ref{thm:arrangement-boundary} and
Section~\ref{sec:features}. Lemma~\ref{lem:grid} constructs the fixed
grid in finitely many rational and finite-field operations.
Lemma~\ref{lem:corner-cone} and finite separation give $D_0$.
Lemma~\ref{lem:edge-wedge} and compact separation give $D_1$.
Lemma~\ref{lem:open-neighborhood} is the crucial bridge to the
strictly positive smooth clearances, and Lemma~\ref{lem:smooth}
gives $D_2$. Section~\ref{sec:effective} makes each choice
computable and finite. The resulting finite octree has exactly the
claimed classes for leaves with nonempty surface intersection. The common interface
construction is finite, and Lemma~\ref{lem:assembly} proves exactness,
embedding, orientation, and feature preservation.
\end{proof}

The dependency chain is, explicitly,
\[
\begin{gathered}
\text{regularized occupancy}
 \Longrightarrow \text{complete S1 complex and genuine strata}\\
\Longrightarrow \text{generic grid and exact corner cones}
 \Longrightarrow \text{frozen corners}\\
\Longrightarrow \text{positive clipped-edge clearances and exact wedges}
 \Longrightarrow \text{closed-contact edge protection}\\
\Longrightarrow \text{open neighborhood of }G
 \Longrightarrow \text{positive planar-remainder clearances}\\
\Longrightarrow \text{finite simple volume cells}
 \Longrightarrow \text{complete common face graphs}\\
\Longrightarrow \text{exact local fillings}
 \Longrightarrow |T|=M.
\end{gathered}
\]

\subsection{Required changes and limits of the construction}

Four changes are indispensable to this proof specification:
\begin{enumerate}
\item protect corner and edge \emph{volume} cells instead of repeatedly
  subdividing all difficult feature neighborhoods;
\item allow arbitrary finite neighbor-depth differences at protected
  boundaries, with no balancing step that reopens protected cells;
\item use complete shared face graphs, including interior loops and
  every required boundary subdivision;
\item replace sampled or MC-style stopping heuristics by exact local
  cone, wedge, or plane certificates.
\end{enumerate}
Exact arrangements supply S1, clearances/certificates, and face traces.
They are not used here to mesh an arbitrary arrangement inside every
leaf: the volume fillings retain only the three proved templates.
This distinguishes Result D from Result C.

The construction deliberately uses uniform depths for each remaining
region and a full finest-depth interface subdivision. These are finite
but can be very expensive. No bound is asserted on practical memory,
running time, mesh size relative to an optimum, minimum triangle angle,
or aspect ratio. In particular exact preservation can retain tiny
segments and produce thin triangles. Any subsequent optimization must
prove that it preserves the established complex and features.

\section{Curved carriers: useful local bounds, still no global theorem}
\label{sec:curved}

The previous sections exploit equality to affine planes, straight
wedges, and exact cones. A sphere, cylinder, general intersection
curve, or smooth blend is not covered merely by substituting its
carrier name into those arguments. The following estimates clarify
some obligations of a future extension; they do not close them.

\subsection{A Hessian bound controls triangle residuals}

\begin{lemma}[Triangle residual bound]
\label{lem:triangle-residual}
Let $\tau=\conv\{p_0,p_1,p_2\}$ and let $f$ be $C^2$ on a
neighborhood of $\tau$. Suppose
\[
  |f(p_i)|\le\eta,\qquad
  \norm{D^2 f(x)}_2\le H\quad(x\in\tau).
\]
Then, for every $x\in\tau$,
\[
  |f(x)|\le \eta+\frac{H\,\diam(\tau)^2}{6}.
  \tag{11}\label{eq:triangle-residual}
\]
\end{lemma}

\begin{proof}
Write $x=\sum_i\lambda_i p_i$ with $\lambda_i\ge0$ and
$\sum_i\lambda_i=1$. Taylor expansion from $x$ to each vertex has
remainder at most $H\norm{p_i-x}^2/2$. The weighted linear terms
cancel, giving
\[
  \left|f(x)-\sum_i\lambda_i f(p_i)\right|
  \le\frac H2\sum_i\lambda_i\norm{p_i-x}^2
  =\frac H2\sum_{i<j}\lambda_i\lambda_j\norm{p_i-p_j}^2.
\]
Since $\sum_{i<j}\lambda_i\lambda_j
=(1-\sum_i\lambda_i^2)/2\le1/3$, the right-hand side is at
most $H\diam(\tau)^2/6$. The weighted vertex residual is at most
$\eta$, proving the claim.
\end{proof}

The bound is for a specified $C^2$ branch with a certified Hessian
bound on the whole triangle, not for a nonsmooth minimum or maximum
across a crease. It controls a field residual; converting that residual
to Euclidean distance needs additional hypotheses.

\subsection{A graph prism converts residual to one-sided distance}

\begin{lemma}[Certified graph-prism estimate]
\label{lem:graph-prism}
Let $e$ be a unit vector, let $D\subset e^\perp$, and consider a
prism of fibers $y+te$ with $y\in D$ and
$\alpha(y)\le t\le\beta(y)$. Suppose $f$ is $C^1$ there and
\[
  \partial_e f\ge m>0,\qquad
  f(y+\alpha(y)e)\le0\le f(y+\beta(y)e)
  \quad\text{for every }y\in D.
\]
Then each fiber has one root $y+g(y)e$. Every prism point $x=y+te$
satisfies
\[
  \dist(x,f^{-1}(0))\le |t-g(y)|\le\frac{|f(x)|}{m}.
  \tag{12}\label{eq:prism-distance}
\]
\end{lemma}

\begin{proof}
The bracket and continuity give existence; strict directional
monotonicity gives uniqueness. Apply the one-variable mean value
theorem on the fiber segment joining $t$ and $g(y)$. That segment
stays in the certified prism, so
$|f(y+te)-f(y+g(y)e)|\ge m|t-g(y)|$.
The first inequality follows because the graph root is a zero of $f$.
\end{proof}

If a triangle lies in such a prism, combining
\eqref{eq:triangle-residual} and~\eqref{eq:prism-distance} gives the
one-sided mesh-to-surface bound
\[
  \sup_{x\in\tau}\dist(x,f^{-1}(0))
  \le \frac{\eta+H\diam(\tau)^2/6}{m}.
\]
The target zero must also be the intended exposed carrier branch.
A residual bound without a certified root bracket and a region on
which the derivative lower bound holds does not justify this step.

Reverse coverage is a separate obligation. For example, if a candidate
mesh is proved to be a graph over \emph{all} of $D$ in the same prism,
and has residual at most $\epsilon$ everywhere, then comparison on
each fiber also bounds the distance from every target graph point to
the mesh by $\epsilon/m$. The residual estimate alone proves neither
that graph property nor coverage of $D$, and does not prevent extra
components, missing features, or a wrong gluing of charts.

\subsection{Pair constraints become ill-conditioned near tangency}

For two unit carrier normals $n_A,n_B$ put $c=n_A\cdot n_B$ and
form the $2\times3$ constraint Jacobian
\[
  J=\begin{pmatrix}n_A^T\\n_B^T\end{pmatrix}.
\]
Since
\[
  JJ^T=\begin{pmatrix}1&c\\c&1\end{pmatrix},
\]
its nonzero singular values are $\sqrt{1+c}$ and $\sqrt{1-c}$;
the smaller is
\[
  \sigma_{\min}(J)=\sqrt{1-|c|}. \tag{13}\label{eq:conditioning}
\]
For two independent affine plane equations, the minimum-norm correction
from residual vector $r$ to their intersection has norm
$\norm{J^+r}\le\norm r/\sigma_{\min}(J)$. The amplification is
unbounded as the planes become parallel or antiparallel. For curved
carriers, the same Jacobian calculation describes local conditioning,
but a finite error theorem additionally needs a certified neighborhood,
control of derivative variation, and existence and uniqueness of the
intended intersection branch. Small pairwise residuals alone are not
a uniform geometric feature-location guarantee.

\subsection{A round-blend radius is not a curvature bound by itself}

For a concrete local example, let $a(x)$ and $b(x)$ be affine signed
plane functions with unit gradients $n_A,n_B$, put $c=n_A\cdot n_B$,
and take the active round-union branch
\[
  f(x)=r-\sqrt{(r-a(x))^2+(r-b(x))^2},\qquad r>0.
\]
This formula applies in the branch where both radial coordinates are
positive. Assume $-1<c<1$ and evaluate at the symmetric zero where
\[
  a=b=r(1-1/\sqrt2).
\]
At that point,
\[
  \nabla f=(n_A+n_B)/\sqrt2,\qquad
  \norm{\nabla f}=\sqrt{1+c},
\]
and direct differentiation gives
\[
  D^2 f=-\frac{(n_A-n_B)(n_A-n_B)^T}{2r}.
\]
The unit tangent across the blend is
$t=(n_A-n_B)/\sqrt{2(1-c)}$. The normal-curvature magnitude in
that direction is therefore
\[
  \kappa=\frac{|t^T(D^2 f)t|}{\norm{\nabla f}}
    =\frac{1-c}{r\sqrt{1+c}}. \tag{14}\label{eq:round-curvature}
\]
The intersection-line direction has zero curvature in this planar
example. Even here $r$ alone does not control curvature uniformly:
as $c\to-1$, the denominator collapses. Other blend definitions or
other branches require their own calculation. This local identity
is not a global curvature certificate for an arbitrary CSG blend.

\subsection{What a curved-SFCC theorem would still need}

A complete extension would have to supply all of the following,
with compatible hypotheses and a termination proof:
\begin{enumerate}
\item complete exposed smooth strata, intersection curves, and corner
  incidence under the intended solid semantics;
\item certified local surface and feature charts, including narrow
  separations and all relevant tangencies or an explicit exclusion
  of them;
\item protected neighborhoods with controlled geometry in place of
  exact cone/wedge equality, and positive bounds that survive
  interface refinement;
\item complete common face traces and error-controlled curve chords,
  with correct ordering and incidence;
\item embedded local patches and relative-boundary homeomorphisms
  or isotopies that agree across every interface;
\item both directions of geometric error control, followed by a
  global gluing and, if claimed, ambient-isotopy argument.
\end{enumerate}
Plantinga--Vegter's smooth regular-zero-set theorem~\cite{pv} does
not by itself supply this composition at nonsmooth CSG junctions.
Nor does the polyhedral protected-cell proof show that its certificates
can simply be replaced by sampled normal variation. The curved extension
remains open in this document.

\section{Relationship to prior work and exact scope}
\label{sec:scope}

The ingredients have substantial prior history. Cubical Marching
Squares~\cite{cms} organizes contouring around shared face segments
and independent cell filling. Its shared-interface idea is relevant
here; it is not used as a substitute for a proof of complete feature
recovery, exact embedding, or termination.

Varadhan et al.~\cite{varadhan} analyze local compatibility conditions
for topology-preserving Marching-Cubes extraction. Those conditions
motivate keeping local topology and interface behavior explicit.
They cannot be applied without their premises, particularly where
feature-aware cells deliberately retain a contour absent from face
corner signs. Plantinga and Vegter~\cite{pv} prove isotopic meshing
for smooth regular implicit surfaces using interval-driven subdivision;
the nonsmooth hard-CSG case here has different local certificates.

Protecting lower-dimensional features before refining higher-dimensional
pieces is established in piecewise-smooth-complex meshing, including
Cheng, Dey, and Ramos~\cite{psc} and Dey and Levine~\cite{dey-levine}.
Their weighted Delaunay constructions are not the octree construction
proved here. They are primary background for protection and termination,
and their distinct hypotheses should not be silently transferred.

Regularized boundary evaluation and rational arrangements are also
established ideas~\cite{cgal-nef}. This document develops a particular
constructive proof for an SFCC-style specification and records its
obstructions. It claims no general absence of prior art and no priority
for the individual ingredients.

\begin{center}
\begin{tabular}{>{\raggedright\arraybackslash}p{0.27\linewidth}
                >{\raggedright\arraybackslash}p{0.64\linewidth}}
\toprule
Statement & Scope\\
\midrule
Orthogonal result & Exact support-coordinate stopping certificate;
  finite balanced refinement; exact simple-cell filling.\\[3pt]
Fixed-grid obstruction & A rational tetrahedron defeats the specified
  face-balanced simple-cell rule at every depth on $[0,1]^3$.\\[3pt]
Arrangement fallback & Exact and finite for any finite octree; permits
  arbitrary local arrangements and coincident interface patches.\\[3pt]
Protected-cell theorem & Exact and finite for rational polyhedral
  manifold boundaries after the explicit protection, interface,
  and certification changes.\\[3pt]
Curved-carrier estimates & Local lemmas and identified obligations;
  no completed curved-SFCC correctness or termination theorem.\\[3pt]
Current implementation & Not certified by this document; no assertion
  about its predicates, feature completeness, finite tolerances,
  or exported floating-point triangles.\\
\bottomrule
\end{tabular}
\end{center}

In particular, rounding the rational mesh to a finite floating-point
format can identify distinct vertices, collapse triangles, or change
an exact intersection. A separate rounding or perturbation theorem
would be needed to preserve these guarantees after conversion.
No mathematical statement above automatically transfers to the app's
GPU sampling or Rust export pipeline.

\appendix
\section{Reproducibility and finite arithmetic cross-checks}

The source is self-contained and uses standard LaTeX packages. From the
repository root, build with generated output outside the repository:
\begin{verbatim}
mkdir -p /tmp/sfcc-correctness-proofs
latexmk -pdf -interaction=nonstopmode -halt-on-error \
  -outdir=/tmp/sfcc-correctness-proofs \
  docs/research/sfcc-correctness-proofs.tex
\end{verbatim}
No application build or runtime rendering is needed for this document.

The orthogonal mask table can be independently reproduced as follows.
Associate a bit with each sign triple in $\{0,1\}^3$. For every pair
of triples differing in one coordinate, add their shared octahedron
edge exactly when their occupancy bits differ. First reject any
octahedron vertex of degree four. Among the remaining nonempty graphs,
reject disconnected graphs. The rest are simple cycles. At a cycle
vertex, a turn occurs exactly when the two neighboring vertices lie
on different coordinate axes. Count cycles by their turns as in
Section~\ref{sec:orthogonal}. Enumeration of all $256$ masks gives
the displayed counts; the table-free argument proves the result.

The tetrahedron checks use exact rational arithmetic: its three base
heights equal $1/6$, its three projected base-line invariants have the
values displayed in Section~\ref{sec:obstruction}, and its apex exit
coordinates have the formula
\[
  (x,y)=(1/3,1/3)+\frac{1/3-z_0}{61}(w_x,w_y),
\]
for $(w_x,w_y)=(0,1),(1,0),(-3,-2)$. Direct rational checks at depths
$2$ through $199$ agree with the same-quarter-face conclusion. Those
finite checks are a transcription check only; inequality
\eqref{eq:tetra-drift} proves the claim for all depths.

\begin{thebibliography}{9}

\bibitem{cms}
C.-C. Ho, F.-C. Wu, B.-Y. Chen, Y.-Y. Chuang, and M. Ouhyoung.
\emph{Cubical Marching Squares: Adaptive Feature Preserving Surface
Extraction from Volume Data.}
Computer Graphics Forum, 24(3), Eurographics 2005.
\url{https://www.csie.ntu.edu.tw/~cyy/publications/papers/Ho2005CMS.pdf}.

\bibitem{cgal-nef}
CGAL Project.
\emph{3D Boolean Operations on Nef Polyhedra: User Manual}, especially
``Regularized Set Operations.''
\url{https://doc.cgal.org/latest/Nef_3/index.html}.
Accessed October 13, 2026.

\bibitem{psc}
S.-W. Cheng, T. K. Dey, and E. A. Ramos.
\emph{Delaunay Refinement for Piecewise Smooth Complexes.}
Author manuscript; conference version, ACM--SIAM Symposium on Discrete
Algorithms (SODA), 2007.
\url{https://cse.hkust.edu.hk/~scheng/pub/soda2007a-psc.pdf}.

\bibitem{dey-levine}
T. K. Dey and J. A. Levine.
\emph{Delaunay Meshing of Piecewise Smooth Complexes without Expensive
Predicates.}
Algorithms, 2(4):1327--1349, 2009.
\url{https://doi.org/10.3390/a2041327}.

\bibitem{varadhan}
G. Varadhan, S. Krishnan, T. V. N. Sriram, and D. Manocha.
\emph{Topology Preserving Surface Extraction Using Adaptive Subdivision.}
Eurographics Symposium on Geometry Processing, 2004.
\url{https://gamma.cs.unc.edu/RECONS/topology.pdf}.

\bibitem{pv}
S. Plantinga and G. Vegter.
\emph{Isotopic Meshing of Implicit Surfaces.}
The Visual Computer, 23:45--58, 2007.
\url{https://doi.org/10.1007/s00371-006-0083-6}.
Author-hosted publication copy:
\url{https://pure.rug.nl/ws/files/2777927/2007VisCompPlantinga.pdf}.

\end{thebibliography}
\end{document}
